\begin{problem}[Characterizations of local rings]
Let $R$ be a commutative ring.  Prove the equivalence of:
\begin{enumerate}
\item $R$ has a unique maximal ideal;
\item the nonunits form an ideal;
\item for every $a\in R$, at least one of $a$ or $1-a$ is a unit.
\end{enumerate}
\end{problem}

\paragraph{Why this problem matters.}
This theorem turns the abstract definition of a local ring into a practical unit test.

\paragraph{Useful ideas before starting.}
Every nonunit belongs to some maximal ideal.  Distinct maximal ideals are comaximal.

\paragraph{Guided hints.}
For the last implication, assume two maximal ideals exist and write $1=a+b$ with one summand in each.

\begin{solution}
If $R$ has unique maximal ideal $\mathfrak m$, every nonunit lies in a maximal ideal and hence in $\mathfrak m$; conversely elements of $\mathfrak m$ cannot be units.  Thus nonunits equal $\mathfrak m$ and form an ideal.  If nonunits form an ideal $N$, then every maximal ideal is contained in $N$ and maximality forces uniqueness.  Finally, locality implies $a$ and $1-a$ cannot both lie in $\mathfrak m$.  Conversely, if the unit alternative holds and $\mathfrak m\ne\mathfrak n$ are maximal, choose $a\in\mathfrak m$, $b\in\mathfrak n$ with $a+b=1$.  Then $a$ and $1-a=b$ are both nonunits, contradiction.
\end{solution}

\paragraph{What happened mathematically.}
The unique maximal ideal is not an extra decoration: it is exactly the set of nonunits.

\paragraph{Extensions.}
Apply the criterion to $\ZZ/n\ZZ$ and characterize exactly when it is local.

\paragraph{Provenance.}
New canonical synthesis problem based on the local-ring and residue-field corpus.
